\section{Citric Acid Cycle (TCA/Krebs Cycle)}
%────────────────────────────────────────────────────────────────

\begin{learningobjectives}
After completing this section, you should be able to:
\begin{enumerate}
    \item Describe the location and overall function of the citric acid cycle
    \item List all eight reactions, enzymes, and products of the TCA cycle
    \item Identify the regulatory enzymes and their allosteric modulators
    \item Calculate the ATP yield from one acetyl-CoA through the TCA cycle
    \item Explain the amphibolic nature of the TCA cycle
    \item Describe anaplerotic reactions that replenish TCA intermediates
\end{enumerate}
\end{learningobjectives}

\subsection{Overview and Location}

\begin{definition}
The \textbf{citric acid cycle} (also called the \textbf{tricarboxylic acid (TCA) cycle} or \textbf{Krebs cycle}) is a series of eight enzymatic reactions that oxidize acetyl-CoA to CO$_2$, generating reduced electron carriers (NADH, FADH$_2$) and GTP.
\end{definition}

\begin{itemize}
    \item \textbf{Location:} Mitochondrial matrix
    \item \textbf{Input:} Acetyl-CoA (2 carbons)
    \item \textbf{Output:} 2 CO$_2$, 3 NADH, 1 FADH$_2$, 1 GTP per acetyl-CoA
    \item \textbf{Function:} Final common pathway for oxidation of carbohydrates, fats, and amino acids
\end{itemize}

\begin{analogy}
The TCA cycle is like a recycling center. Oxaloacetate (OAA) is the ``container'' that picks up acetyl-CoA (the ``trash''). As the container moves through the cycle, the trash is processed (oxidized), energy is extracted (NADH, FADH$_2$, GTP), and waste products are released (CO$_2$). At the end, the empty container (OAA) is ready to pick up more trash!
\end{analogy}

\subsection{Pyruvate Dehydrogenase Complex (PDC)}

Before entering the TCA cycle, pyruvate must be converted to acetyl-CoA.

\begin{theorem}
\textbf{Pyruvate Dehydrogenase Reaction:}
\[
\text{Pyruvate} + \text{CoA} + \text{NAD}^+ \xrightarrow{\text{PDC}} \text{Acetyl-CoA} + \text{CO}_2 + \text{NADH}
\]

\begin{itemize}
    \item $\Delta G°' = -33.4$ kJ/mol (irreversible)
    \item Location: Mitochondrial matrix
    \item Links glycolysis to TCA cycle
\end{itemize}
\end{theorem}

\subsubsection{PDC Structure and Mechanism}

The PDC is a massive multi-enzyme complex with three catalytic enzymes:

\begin{center}
\begin{tabular}{clll}
\toprule
\textbf{Component} & \textbf{Enzyme} & \textbf{Coenzyme} & \textbf{Vitamin Source} \\
\midrule
E1 & Pyruvate dehydrogenase & TPP & Thiamine (B$_1$) \\
E2 & Dihydrolipoyl transacetylase & Lipoic acid, CoA & Pantothenic acid (B$_5$) \\
E3 & Dihydrolipoyl dehydrogenase & FAD, NAD$^+$ & Riboflavin (B$_2$), Niacin (B$_3$) \\
\bottomrule
\end{tabular}
\end{center}

\begin{tcolorbox}[colback=yellow!10, colframe=yellow!70!black, title=\textbf{Memory Aid: PDC Coenzymes}]
\textbf{``Tender Loving Care For Nancy''}
\begin{itemize}
    \item \textbf{T}PP (Thiamine)
    \item \textbf{L}ipoic acid
    \item \textbf{C}oA (pantothenic acid)
    \item \textbf{F}AD (riboflavin)
    \item \textbf{N}AD$^+$ (niacin)
\end{itemize}
\end{tcolorbox}

\subsubsection{Regulation of PDC}

\begin{theorem}
\textbf{PDC Regulation by Covalent Modification:}

\begin{center}
\begin{tabular}{ll}
\toprule
\textbf{Active (dephosphorylated)} & \textbf{Inactive (phosphorylated)} \\
\midrule
PDC phosphatase activated by: & PDC kinase activated by: \\
\quad Ca$^{2+}$ (muscle contraction) & \quad ATP (high energy) \\
\quad Insulin & \quad Acetyl-CoA (product) \\
\quad Pyruvate (substrate) & \quad NADH (product) \\
& \quad Fatty acids (alternative fuel) \\
\bottomrule
\end{tabular}
\end{center}
\end{theorem}

\begin{clinmark}
\textbf{Clinical Correlation: Pyruvate Dehydrogenase Deficiency}

PDC deficiency is an X-linked disorder causing:
\begin{itemize}
    \item Lactic acidosis (pyruvate converted to lactate)
    \item Neurological defects (brain depends on glucose oxidation)
    \item Treatment: Ketogenic diet (provides acetyl-CoA from fats, bypassing PDC)
\end{itemize}
\end{clinmark}

\subsection{Reactions of the TCA Cycle}

\begin{center}
\begin{tcolorbox}[colback=green!5, colframe=green!70!black, title=\textbf{TCA Cycle Overview}]
\textbf{Net Reaction (per acetyl-CoA):}
\[
\text{Acetyl-CoA} + 3\text{NAD}^+ + \text{FAD} + \text{GDP} + \text{P}_i + 2\text{H}_2\text{O}
\]
\[
\downarrow
\]
\[
2\text{CO}_2 + 3\text{NADH} + \text{FADH}_2 + \text{GTP} + \text{CoA} + 3\text{H}^+
\]
\end{tcolorbox}
\end{center}

\subsubsection{The Eight Reactions}

\textbf{Reaction 1: Citrate Synthase}

\[
\text{Oxaloacetate} + \text{Acetyl-CoA} + \text{H}_2\text{O} \rightarrow \text{Citrate} + \text{CoA-SH}
\]

\begin{itemize}
    \item $\Delta G°' = -31.4$ kJ/mol (\textbf{irreversible})
    \item \textbf{Type:} Condensation (aldol condensation)
    \item 4C + 2C $\rightarrow$ 6C
    \item \textbf{Regulatory enzyme:} Inhibited by ATP, NADH, succinyl-CoA, citrate
\end{itemize}

\textbf{Reaction 2: Aconitase}

\[
\text{Citrate} \xleftrightarrow{\text{Aconitase}} \text{Isocitrate}
\]

\begin{itemize}
    \item $\Delta G°' = +6.3$ kJ/mol (reversible)
    \item \textbf{Type:} Isomerization (via cis-aconitate intermediate)
    \item Moves hydroxyl group from C3 to C2
    \item Contains iron-sulfur cluster (Fe-S)
\end{itemize}

\begin{clinmark}
\textbf{Clinical Note:} Fluoroacetate (rat poison) is converted to fluorocitrate, which inhibits aconitase, blocking the TCA cycle.
\end{clinmark}

\textbf{Reaction 3: Isocitrate Dehydrogenase (IDH)}

\[
\text{Isocitrate} + \text{NAD}^+ \rightarrow \alpha\text{-Ketoglutarate} + \text{CO}_2 + \text{NADH}
\]

\begin{itemize}
    \item $\Delta G°' = -20.9$ kJ/mol (\textbf{irreversible})
    \item \textbf{Type:} Oxidative decarboxylation
    \item 6C $\rightarrow$ 5C (first CO$_2$ released)
    \item \textbf{First NADH produced}
    \item \textbf{Rate-limiting step of TCA cycle}
    \item \textbf{Regulation:}
    \begin{itemize}
        \item Activated by: ADP, Ca$^{2+}$
        \item Inhibited by: ATP, NADH
    \end{itemize}
\end{itemize}

\textbf{Reaction 4: $\alpha$-Ketoglutarate Dehydrogenase Complex}

\[
\alpha\text{-Ketoglutarate} + \text{NAD}^+ + \text{CoA} \rightarrow \text{Succinyl-CoA} + \text{CO}_2 + \text{NADH}
\]

\begin{itemize}
    \item $\Delta G°' = -33.5$ kJ/mol (\textbf{irreversible})
    \item \textbf{Type:} Oxidative decarboxylation
    \item 5C $\rightarrow$ 4C (second CO$_2$ released)
    \item \textbf{Second NADH produced}
    \item Same coenzymes as PDC: TPP, lipoic acid, CoA, FAD, NAD$^+$
    \item \textbf{Regulation:}
    \begin{itemize}
        \item Activated by: Ca$^{2+}$
        \item Inhibited by: Succinyl-CoA, NADH, ATP
    \end{itemize}
\end{itemize}

\textbf{Reaction 5: Succinyl-CoA Synthetase}

\[
\text{Succinyl-CoA} + \text{GDP} + \text{P}_i \xleftrightarrow{} \text{Succinate} + \text{CoA} + \text{GTP}
\]

\begin{itemize}
    \item $\Delta G°' = -2.9$ kJ/mol (reversible)
    \item \textbf{Type:} Substrate-level phosphorylation
    \item \textbf{Only GTP (or ATP) produced directly in TCA cycle}
    \item High-energy thioester bond drives GTP synthesis
    \item GTP can be converted to ATP: GTP + ADP $\rightleftharpoons$ GDP + ATP
\end{itemize}

\textbf{Reaction 6: Succinate Dehydrogenase (Complex II)}

\[
\text{Succinate} + \text{FAD} \xleftrightarrow{} \text{Fumarate} + \text{FADH}_2
\]

\begin{itemize}
    \item $\Delta G°' = 0$ kJ/mol (reversible)
    \item \textbf{Type:} Oxidation (dehydrogenation)
    \item \textbf{Only FADH$_2$ produced in TCA cycle}
    \item \textbf{Only enzyme embedded in inner mitochondrial membrane}
    \item Part of electron transport chain (Complex II)
    \item FAD used because reaction not energetic enough to reduce NAD$^+$
    \item \textbf{Competitively inhibited by malonate} (succinate analog)
\end{itemize}

\textbf{Reaction 7: Fumarase}

\[
\text{Fumarate} + \text{H}_2\text{O} \xleftrightarrow{} \text{L-Malate}
\]

\begin{itemize}
    \item $\Delta G°' = -3.8$ kJ/mol (reversible)
    \item \textbf{Type:} Hydration
    \item Stereospecific: produces only L-malate
\end{itemize}

\textbf{Reaction 8: Malate Dehydrogenase}

\[
\text{L-Malate} + \text{NAD}^+ \xleftrightarrow{} \text{Oxaloacetate} + \text{NADH} + \text{H}^+
\]

\begin{itemize}
    \item $\Delta G°' = +29.7$ kJ/mol (highly unfavorable!)
    \item \textbf{Type:} Oxidation
    \item \textbf{Third NADH produced}
    \item Reaction pulled forward by:
    \begin{itemize}
        \item Low [OAA] — rapidly consumed by citrate synthase
        \item Removal of NADH by ETC
    \end{itemize}
    \item Regenerates oxaloacetate to continue the cycle
\end{itemize}

\begin{tcolorbox}[colback=blue!5, colframe=blue!70!black, title=\textbf{Memory Aid: TCA Cycle Intermediates}]
\textbf{``Citrate Is Krebs' Starting Substrate For Making Oxaloacetate''}

\begin{itemize}
    \item \textbf{C}itrate (6C)
    \item \textbf{I}socitrate (6C)
    \item $\alpha$-\textbf{K}etoglutarate (5C)
    \item \textbf{S}uccinyl-CoA (4C)
    \item \textbf{S}uccinate (4C)
    \item \textbf{F}umarate (4C)
    \item \textbf{M}alate (4C)
    \item \textbf{O}xaloacetate (4C)
\end{itemize}
\end{tcolorbox}

\subsection{Energy Yield of the TCA Cycle}

\begin{center}
\begin{tabular}{lcc}
\toprule
\textbf{Product} & \textbf{Per Acetyl-CoA} & \textbf{ATP Equivalent} \\
\midrule
NADH & 3 & 3 $\times$ 2.5 = 7.5 ATP \\
FADH$_2$ & 1 & 1 $\times$ 1.5 = 1.5 ATP \\
GTP & 1 & 1 ATP \\
\midrule
\textbf{Total} & & \textbf{10 ATP} \\
\bottomrule
\end{tabular}
\end{center}

\textbf{Per glucose molecule (2 acetyl-CoA):}
\begin{itemize}
    \item 6 NADH $\rightarrow$ 15 ATP
    \item 2 FADH$_2$ $\rightarrow$ 3 ATP
    \item 2 GTP $\rightarrow$ 2 ATP
    \item \textbf{Total from TCA: 20 ATP per glucose}
\end{itemize}

\subsection{Regulation of the TCA Cycle}

\begin{theorem}
\textbf{Three Regulatory Enzymes:}

\begin{enumerate}
    \item \textbf{Citrate synthase} (entry point)
    \item \textbf{Isocitrate dehydrogenase} (rate-limiting)
    \item \textbf{$\alpha$-Ketoglutarate dehydrogenase}
\end{enumerate}

\textbf{General
principles:}
\begin{itemize}
    \item High ATP, NADH, succinyl-CoA, citrate $\rightarrow$ \textbf{inhibit} cycle
    \item High ADP, NAD$^+$, Ca$^{2+}$ $\rightarrow$ \textbf{activate} cycle
\end{itemize}
\end{theorem}

\begin{center}
\begin{tabular}{lccc}
\toprule
\textbf{Regulator} & \textbf{Citrate Synthase} & \textbf{Isocitrate DH} & \textbf{$\alpha$-KG DH} \\
\midrule
ATP & Inhibits & Inhibits & — \\
ADP & — & Activates & — \\
NADH & Inhibits & Inhibits & Inhibits \\
NAD$^+$ & — & Activates & Activates \\
Succinyl-CoA & — & — & Inhibits \\
Ca$^{2+}$ & — & Activates & Activates \\
Citrate & Inhibits & — & — \\
\bottomrule
\end{tabular}
\end{center}

\begin{analogy}
The TCA cycle is like a factory production line. When the warehouse is full of products (high ATP, NADH), the line slows down. When orders are piling up (high ADP, NAD$^+$), production speeds up. Calcium is like an ``urgent order'' signal—during muscle contraction, Ca$^{2+}$ increases, signaling the need for more ATP.
\end{analogy}

\subsection{Anaplerotic Reactions}

\begin{definition}
\textbf{Anaplerotic reactions} are reactions that replenish TCA cycle intermediates when they are removed for biosynthesis. ``Anaplerotic'' means ``filling up.''
\end{definition}

\textbf{Key Anaplerotic Reactions:}

\begin{enumerate}
    \item \textbf{Pyruvate carboxylase} (most important):
    \[
    \text{Pyruvate} + \text{CO}_2 + \text{ATP} \xrightarrow{\text{Biotin}} \text{Oxaloacetate} + \text{ADP} + \text{P}_i
    \]
    \begin{itemize}
        \item Activated by acetyl-CoA
        \item Located in mitochondria
        \item Also used in gluconeogenesis
    \end{itemize}
    
    \item \textbf{PEP carboxykinase} (in some tissues):
    \[
    \text{PEP} + \text{CO}_2 \rightarrow \text{Oxaloacetate} + \text{P}_i
    \]
    
    \item \textbf{Amino acid degradation}:
    \begin{itemize}
        \item Glutamate $\rightarrow$ $\alpha$-ketoglutarate
        \item Aspartate $\rightarrow$ oxaloacetate
    \end{itemize}
\end{enumerate}

\subsection{TCA Cycle as a Metabolic Hub}

\begin{theorem}
The TCA cycle intermediates serve as precursors for many biosynthetic pathways:

\begin{itemize}
    \item \textbf{Citrate} $\rightarrow$ Fatty acid synthesis (exported to cytoplasm)
    \item \textbf{$\alpha$-Ketoglutarate} $\rightarrow$ Glutamate $\rightarrow$ other amino acids, purines
    \item \textbf{Succinyl-CoA} $\rightarrow$ Heme synthesis, odd-chain fatty acid metabolism
    \item \textbf{Oxaloacetate} $\rightarrow$ Aspartate $\rightarrow$ pyrimidines, other amino acids
    \item \textbf{Oxaloacetate} $\rightarrow$ Gluconeogenesis (via PEP)
\end{itemize}
\end{theorem}

\begin{tcolorbox}[colback=green!5, colframe=green!70!black, title=\textbf{Clinical Correlation: TCA Cycle Defects}]
\textbf{Mutations in TCA cycle enzymes} are associated with various cancers:
\begin{itemize}
    \item \textbf{Succinate dehydrogenase (SDH) mutations}: Paragangliomas, pheochromocytomas
    \item \textbf{Fumarase mutations}: Renal cell carcinoma, leiomyomas
    \item \textbf{Isocitrate dehydrogenase (IDH) mutations}: Gliomas, acute myeloid leukemia
\end{itemize}
These mutations lead to accumulation of ``oncometabolites'' that promote tumor growth.
\end{tcolorbox}

\subsection{Practice Problems: TCA Cycle}

\begin{enumerate}
    \item A patient has a thiamine deficiency. Which TCA cycle reactions would be directly affected?
    
    \item Calculate the total ATP yield from one molecule of acetyl-CoA through the TCA cycle and oxidative phosphorylation.
    
    \item Why does the TCA cycle stop when oxygen is absent, even though oxygen doesn't directly participate in any TCA cycle reaction?
    
    \item A researcher adds fluoroacetate (converted to fluorocitrate, an aconitase inhibitor) to cells. What would happen to:
    \begin{enumerate}
        \item Citrate levels?
        \item Isocitrate levels?
        \item ATP production?
    \end{enumerate}
\end{enumerate}

\textbf{Solutions:}
\begin{enumerate}
    \item Thiamine (vitamin B$_1$) is a component of TPP, required for:
    \begin{itemize}
        \item Pyruvate dehydrogenase (before TCA)
        \item $\alpha$-Ketoglutarate dehydrogenase (in TCA)
    \end{itemize}
    Both reactions would be impaired.
    
    \item Per acetyl-CoA:
    \begin{itemize}
        \item 3 NADH $\times$ 2.5 ATP = 7.5 ATP
        \item 1 FADH$_2$ $\times$ 1.5 ATP = 1.5 ATP
        \item 1 GTP = 1 ATP
        \item \textbf{Total: 10 ATP}
    \end{itemize}
    
    \item Without oxygen, the electron transport chain cannot function, so NADH and FADH$_2$ cannot be reoxidized to NAD$^+$ and FAD. Without NAD$^+$ and FAD, the dehydrogenase reactions of the TCA cycle cannot proceed.
    
    \item \begin{enumerate}
        \item Citrate levels would \textbf{increase} (accumulate before block)
        \item Isocitrate levels would \textbf{decrease} (cannot be formed)
        \item ATP production would \textbf{decrease dramatically} (cycle blocked)
    \end{enumerate}
\end{enumerate}

